% AIPL に express mode と where を入れるとどうなるか
%   lualatex -interaction=nonstopmode aipl_abcl1_features_ja.tex   （2回）
\documentclass[11pt,a4paper]{ltjsarticle}
\usepackage{amsmath,amssymb}
\usepackage{amsthm}
\usepackage{mathpartir}
\usepackage{listings}
\usepackage{xcolor}
\usepackage{geometry}
\usepackage{array}
\usepackage{booktabs}
\usepackage[hidelinks]{hyperref}
\geometry{margin=22mm}

\theoremstyle{remark}
\newtheorem{remark}{注意}
\newtheorem{finding}{測定}

\definecolor{codebg}{gray}{0.96}
\lstset{
  basicstyle=\ttfamily\small,
  backgroundcolor=\color{codebg},
  frame=single,
  framesep=4pt,
  rulecolor=\color{gray},
  breaklines=true,
  showstringspaces=false,
  columns=fullflexible,
  keepspaces=true,
  keywordstyle=\bfseries,
  commentstyle=\itshape\color{gray!130},
}
\lstdefinelanguage{AIPL}{
  morekeywords={class,method,var,float,new,now,future,await,reply,send,become,
    call,if,else,while,do,select,case,timeout,self,sender,remote,print,
    where,express,atomic,int,string,bool,unit,any},
  sensitive=true,
  morecomment=[l]{//},
  morestring=[b]",
}

\newcommand{\ty}[1]{\ensuremath{\mathtt{#1}}}

\title{AIPL の言語仕様の検討\\
{\large ---\ ABCL/1 の機構の再導入と、分散 AI エージェント／実時間 OS という目的から ---}}
\author{ABCL$^{c}$+ / AIOS 処理系（\texttt{yaskodama/abclcp}）}
\date{2026年07月30日 09:48 JST 版}

\begin{document}
\maketitle

\begin{abstract}
先の比較レポートで、AIPL は ABCL/1 から三種のメッセージ送信を継承しつつ
express mode、返信先の一級性、\texttt{where} 制約つきパターン受信を
落としていることを整理した。本稿はそのうち
\textbf{express mode}と\textbf{\texttt{where}}の二つを AIPL に入れた場合に
何が起きるかを、構文・実装・型システム・形式的保証の四層で検討する。

結論を三点にまとめる。
第一に、\textbf{両者は型システムをほとんど壊さない}。どちらも
「いつメッセージを受理するか」の話であり、「どんな型を持つか」の話ではないためである。
第二に、\texttt{where} は\textbf{進行定理に第四の枝「ガード待ち」を追加する}。
これは失うものが明確で、しかも \texttt{select} と名前 dispatch の並立という
既存の不整合を解消できるので、費用対効果が高い。
第三に、express mode は意外にも\textbf{型健全性の証明を壊さない}。
形式モデル AIPL$^-$ はアクターの逐次性を最初から落としており、
同一アクターへの複数タスクの並行実行を既に許しているからである。
壊れるのは\emph{プログラム側の}不変条件であり、だからこそ ABCL/1 は
atomicity 宣言を用意していた。

一方で express mode には実装上の前提がある。現在のメソッドのローカル変数は
アクターの環境に直接置かれ、呼び出し時に退避・復元されていない。
これは今は観測できない潜在的な問題だが、入れ子実行を導入した途端に
表に出る。本稿はこれを実測で確認した。

第 \ref{part:purpose} 部では視点を変え、AIPL が
\textbf{分散 AI エージェント言語}と
\textbf{実時間 OS（フィジカル AI）の記述言語}という
二つの目的を持つとしたら、どのような言語仕様であるべきかを考察する。
結論は\textbf{一つの中核と二つのプロファイル}である。
二つの目的の要求は多くの点で正反対であり、
一つの言語で両方を無差別に許すと、どちらの保証も得られない。
両者を静的に分ける鍵は effect/capability 系であり、
幸いなことに実装にはすでに 17 種の capability 分類が存在する
（ただし現在は実行時メタデータにとどまり、型検査器は使っていない）。
Ada の Ravenscar プロファイルという直接の先例もある。
\end{abstract}

\tableofcontents

\part{express mode と \texttt{where} の再導入}
\label{part:mech}

\section{はじめに}

\subsection{問い}

AIPL（ABCL$^{c}$+）は ABCL/1 の子孫を自認しているが、
ABCL/1 の中核機構のうち次の三つを持っていない。

\begin{itemize}
\item \textbf{express mode} ---- 処理中のメッセージを中断して割り込む機構。
\item \textbf{返信先の一級性} ---- 返信先を値として他者へ渡し、委譲する機構。
\item \textbf{\texttt{where}} ---- オブジェクトの状態に依存した受理条件。
\end{itemize}

本稿は前者と後者、すなわち\emph{受理と割り込みに関する二つ}を導入した場合の
帰結を検討する。返信先の一級化は線形型を要する別問題なので扱わない。

\subsection{なぜこの二つなのか}

現在の AIPL には次の具体的な穴がある。

\begin{itemize}
\item \texttt{now} で返信を待っているアクターは、その間\textbf{他のメッセージを
  一切処理できない}。1 アクター 1 スレッドで、\texttt{await\_future} が
  条件変数でブロックするためである。状態を問い合わせることも、
  止めることもできない。
\item 状態に依存した待ち合わせが書けない。「バッファが空でないときだけ
  \texttt{get} を受理する」を表現するには、受理してから条件を見て、
  駄目なら自分に送り直すしかない。受信規律の外へ出てしまう。
\item 受信経路が二本ある。名前による dispatch と \texttt{select} が並立し、
  同名のメッセージが通常の dispatch に先取りされると \texttt{select} には
  残らない。
\end{itemize}

express mode は一つ目、\texttt{where} は二つ目と三つ目に直接効く。

\section{提案する形}

\subsection{\texttt{where} ---- 状態依存の受理条件}

メソッド宣言と \texttt{select} の case に、真偽値の\emph{ガード}を付けられる
ようにする。

\begin{lstlisting}[language=AIPL]
class Buffer {
  var n = 0;
  var cap = 3;

  // n < cap のときだけ put を受理する
  method put(x) : unit where n < cap {
    n = n + 1;
    print("put, n=" ++ n);
  }

  // n > 0 のときだけ get を受理する
  method get() : int where n > 0 {
    n = n - 1;
    reply(n);
  }
}
\end{lstlisting}

\texttt{select} 側も同じ形にする。

\begin{lstlisting}[language=AIPL]
select {
  case get() where n > 0 -> { reply(n); }
  timeout 500 -> { print("empty"); }
}
\end{lstlisting}

意味は ABCL/1 と同じである。ガードが偽のメッセージは\textbf{受理されず、
メールボックスに残る}。状態が変わって真になれば受理される。

\subsection{express mode ---- 割り込み}

送信側は三種それぞれに express 版を用意する。既存の \texttt{send!}
（\texttt{UnsafeSend}）と衝突しないよう、修飾子として \texttt{express} を置く。

\begin{lstlisting}[language=AIPL]
express send a.stop();          // 割り込んで実行、待たない
var s = express now a.status(); // 割り込んで実行、返信を待つ
var f = express future a.status();
\end{lstlisting}

受信側は、割り込みで実行されてよいメソッドを明示する。
どのメソッドが他のメソッドの実行を中断しうるかは、
読み手にとって最も重要な情報だからである。

\begin{lstlisting}[language=AIPL]
class Worker {
  var progress = 0;
  var stopped  = 0;

  // 通常のメソッド。長く走る
  method run(n) : unit {
    var i = 0;
    while i < n do {
      if (stopped >= 1) { i = n; }
      progress = progress + 1;
      i = i + 1;
    }
  }

  // 割り込みで呼ばれてよい。run の実行中でも応答する
  express method status() : int {
    reply(progress);
  }

  express method stop() : unit {
    stopped = 1;
  }
}
\end{lstlisting}

割り込まれたくない区間は \texttt{atomic} で保護する。

\begin{lstlisting}[language=AIPL]
method transfer(k) : unit {
  atomic {
    from = from - k;      // この2行の間に割り込まれると
    to   = to + k;        // 合計が壊れる
  }
}
\end{lstlisting}

\section{実装への影響}

AIPL の実装を読んだうえで、変更箇所を具体的に見積もる。

\subsection{\texttt{where} の実装}

\begin{center}
\small
\begin{tabular}{@{}lp{9.4cm}@{}}
\toprule
ファイル & 変更 \\
\midrule
\texttt{lexer.mll}   & \texttt{where} をキーワードに追加 \\
\texttt{parser.mly}  & \texttt{method\_decl} に省略可能な \texttt{where expr}、
  \texttt{select\_case} にも同じものを追加 \\
\texttt{ast.ml}      & \texttt{method\_decl} に \texttt{guard : expr option}、
  \texttt{select\_case} にも同様 \\
\texttt{infer.ml}    & ガードをメソッドの環境（フィールド＋仮引数）で
  \ty{bool} として型検査するだけ \\
\texttt{eval\_thread.ml} & \textbf{ここが本体}。\texttt{actor\_loop} が
  \texttt{Queue.pop} で先頭を取るのをやめ、メールボックスを走査して
  「名前が一致し、かつガードが真」の最初のメッセージを選ぶ \\
\bottomrule
\end{tabular}
\end{center}

走査の機構はすでにある。\texttt{select} の実装 \texttt{pop\_matching\_message}
はメールボックスを一旦すべて取り出し、パターンに合うものを探し、
残りを戻している。ガード評価を足すだけで再利用できる。

\begin{remark}
これは \texttt{select} と通常 dispatch を\textbf{統一する好機}でもある。
通常 dispatch がガード付きの走査になるなら、\texttt{select} は
「この名前の集合のみを対象に、タイムアウト付きで走査する」という
特別な場合にすぎない。受信経路が一本になり、
先取り競合という現在の不整合が消える。
\end{remark}

\subsection{express mode の実装}

\begin{center}
\small
\begin{tabular}{@{}lp{9.4cm}@{}}
\toprule
ファイル & 変更 \\
\midrule
\texttt{lexer.mll}   & \texttt{express}, \texttt{atomic} をキーワードに追加 \\
\texttt{parser.mly}  & 送信の express 版 3 種、\texttt{express method}、
  \texttt{atomic \{ ... \}} 文 \\
\texttt{ast.ml}      & 送信に mode を持たせる、\texttt{method\_decl} に
  \texttt{express : bool}、\texttt{Atomic of stmt} \\
\texttt{types.ml} / \texttt{infer.ml} &
  express メソッドの \texttt{reply} を、選択されたメッセージの $\rho$ に
  束縛する（\texttt{select} の case で既に行っている処理と同型） \\
\texttt{eval\_thread.ml} &
  アクターに第二のキューを足す。\texttt{eval\_stmt} の文境界で
  express キューを覗き、あれば\emph{入れ子で}評価する。
  \texttt{atomic} 区間ではこの覗き込みを止める \\
\bottomrule
\end{tabular}
\end{center}

割り込みは協調的なポーリングで実現する。OCaml のシステムスレッドを
外から割り込むのは現実的でないので、\texttt{eval\_stmt} が文の境界で
express キューを確認し、あればその場で express メソッドの本体を
入れ子に評価して戻る。粒度は文単位になる。

\subsection{前提となる修正 ---- ローカル変数のスコープ}

入れ子実行を入れる前に直さなければならない箇所がある。

\begin{finding}
現在、メソッドの\textbf{仮引数}は呼び出しの前後で退避・復元されている。

\begin{lstlisting}
let saved = List.map (fun p -> (p, Hashtbl.find_opt actor.env p)) params in
  List.iter2 (fun p v -> Hashtbl.replace actor.env p v) params arg_vals;
  ...
  eval_stmt actor mdecl.body;
  List.iter (fun (p, ov) -> ...) saved      (* 復元 *)
\end{lstlisting}

しかし本体の \texttt{var} で宣言された\textbf{ローカル変数}は
\texttt{eval\_stmt} が \texttt{actor.env} へ直接書き込むだけで、
退避も復元もされない。
\end{finding}

これは\emph{今は観測できない}。型検査器が各メソッドを独立した環境
（\texttt{clone env}）で検査するため、他のメソッドのローカルを参照する
プログラムは型エラーになる。実測すると次のとおりである。

\begin{lstlisting}[language=AIPL]
class C {
  method a() : unit { var tmp = 777; print("a: tmp=" ++ tmp); }
  method b() : unit { print("b: tmp=" ++ tmp); }   // 宣言していない
}
\end{lstlisting}

\begin{lstlisting}[language={}]
TYPE_ERROR: line 8, col 24: unbound variable: tmp
\end{lstlisting}

型が守っているので実害が出ていない。しかし express mode を入れると
\textbf{同一アクター上で入れ子実行が起きる}。
外側のメソッドが \texttt{var tmp} を持ち、割り込んだ express メソッドも
\texttt{var tmp} を持てば、両者は型検査上まったく別の変数なのに
実行時には同じ \texttt{actor.env} のエントリを共有し、上書きされる。

したがって express mode の前提として、
\textbf{メソッド呼び出しごとのフレームを導入する}
（またはローカル変数も仮引数と同様に退避・復元する）修正が必要である。
これは express mode とは独立に価値のある修正でもある。

\section{型システムへの影響 ---- ほぼ無傷}

ここが本稿の中心的な観察である。

\subsection{どちらも「型」の話ではない}

\texttt{where} も express mode も、\textbf{メッセージをいつ受理するか}を
変える機構であって、\textbf{メッセージがどんな型を持つか}を変えない。
AIPL の型システムが述べているのは後者だけである。

\begin{center}
\small
\begin{tabular}{@{}lp{8.4cm}@{}}
\toprule
現在の検査 & 二機構の影響 \\
\midrule
メソッドの存在       & 影響なし \\
引数の個数と型       & 影響なし \\
返信の型             & 影響なし \\
返信の回数（高々一度）& express は\emph{文脈の分離}が必要。
  \texttt{select} の case で既に実装済みの処理と同型 \\
\texttt{select} の照合 & \texttt{where} でガードの \ty{bool} 検査が増えるだけ \\
公開時の注釈         & 影響なし \\
\bottomrule
\end{tabular}
\end{center}

\subsection{\texttt{where} の型付け}

ガードはメソッドの環境（フィールド＋仮引数）で \ty{bool} として
型検査するだけである。追加規則は次の一行で足りる。

\begin{mathpar}
\inferrule*[right=\textsc{Method-Guard}]
  {\Gamma, \overline{x{:}\tau} \vdash g : \ty{bool} \\
   \Gamma, \overline{x{:}\tau} \vdash_{C.m} \mathit{body}}
  {\Gamma \vdash \texttt{method}\ m(\overline{x})\ \texttt{where}\ g\ \{\mathit{body}\}}
\end{mathpar}

依存型や篩型は\emph{不要}である。それらが必要になるのは
「ガードが真なら本体は安全」を\emph{証明}したい場合であって、
型安全性のためではない。

\subsection{express mode における \texttt{reply} の帰属}

express メソッドの本体にある \texttt{reply} は、割り込んだ express
メッセージへの返信である。囲んでいた通常メソッドの返信ではない。

これは既に解決済みの問題と同型である。\texttt{select} の case 本体の
\texttt{reply} も、囲むメソッドではなく選択されたメッセージへの返信であり、
型検査器は case 本体で \texttt{reply} をそのメソッドの $\rho$ に
束縛し直している。express メソッドも同じ扱いでよい。
実行時に \texttt{ctx\_msg\_id} を退避・復元する処理も、
\texttt{select} の実装がそのまま雛形になる。

\subsection{ガードの純粋性という新しい要求}

一点だけ、型システムに新しい要求が生じる。
\textbf{ガードは副作用を持ってはならない}。
受理の可否を試すためにガードを評価するので、
そこで状態が変わると「試したら受理条件が変わった」ことになる。

現在の OCaml 実装には効果システムが無い。最小限の対処は構文的な制限である。

\begin{itemize}
\item ガードで参照できるのはフィールドと仮引数のみ。
\item 呼び出せるのは純粋なプリミティブのみ（比較、算術など）。
\item \texttt{now} / \texttt{future} / \texttt{send} / \texttt{reply} /
  \texttt{new} / \texttt{become} は書けない。
\end{itemize}

これは既存の構文検査（\texttt{expr\_has\_reply} など）と同じ形の
走査で実装できる。

\section{形式的保証への影響}

\subsection{保存定理 ---- 影響なし}

保存（preservation）は「well-typed な構成が一歩進んでも well-typed」である。
状態の型不変条件は「状態が常に宣言された型を持つ」ことしか言わない。
\texttt{where} は遷移が起こる条件を狭めるだけ、
express mode は遷移の順序を変えるだけなので、どちらも保存を壊さない。

\subsection{進行定理 ---- \texttt{where} は枝を一本増やす}

現在の進行定理は三択として述べられている ---- 値・簡約・\texttt{await} 待ち。
デッドロックは型では排除できないので、その境界を明示したものである。

\texttt{where} を入れると、\textbf{第四の枝が必要になる}。
メールボックスにメッセージがあるのに、どのガードも真でないため
配送できない構成が存在するからである。

\begin{center}
\begin{tabular}{@{}ll@{}}
\toprule
現在 & \texttt{where} 導入後 \\
\midrule
値である & 値である \\
簡約できる & 簡約できる \\
\texttt{await} 待ちである & \texttt{await} 待ちである \\
 & \textbf{ガード待ちである} \\
\bottomrule
\end{tabular}
\end{center}

これは失うものが明確である。定理の主張が一つ弱くなる代わりに、
「なぜ止まっているか」の分類が一つ増える。
言い換えれば、\texttt{where} は\textbf{デッドロックの形を一つ増やす}。
バッファが空のまま \texttt{get} だけが溜まれば、それは正当な待ちであるが、
プログラムの誤りでそうなった場合と静的に区別できない。

\subsection{逐次性 ---- express mode は証明を壊さない}

これは直観に反するが、確認できた事実である。
形式モデル AIPL$^-$ は「保証しないこと」の一覧で、
アクターの逐次性を明示的に\emph{落として}いる。

\begin{quote}
\small
actor の逐次性 ---- 落とした。同じ actor あての二通が同時にタスクになりうる。
\end{quote}

配送規則は無条件に新しいタスクを起こすので、
同一アクターに対して複数のタスクが並行に走りうる。
そして同レポートは、これが\emph{型安全性には影響しない}と明言している。
状態の型不変条件は状態が常に正しい型であることしか言わず、
どの順で書き換えられるかは問わないからである。

したがって express mode の割り込みは、
\textbf{形式モデルが既に許している実行の一部}にすぎない。
逆に言えば、現在の\emph{実装}が形式モデルより厳しく
（1 アクター 1 スレッドの逐次実行）作られているだけである。
express mode はその差を埋める方向の変更であり、
型健全性の証明をやり直す必要はない。

\subsection{危ういのはプログラム側の不変条件}

壊れるのは型ではなく、プログラムが持っている\emph{意味的な}不変条件である。

リポジトリの健全性レポートは哲学者の食事問題について資源デッドロック自由を
証明しているが、その証明はフォーク割り当てプロトコルの抽象状態機械
$s = (\mathrm{phase}, \mathrm{own})$ についてのものであり、
状態の健全性 $\mathrm{wf}(s)$ は
「\textsf{Hold} なら \texttt{lo} のみを持つ」のように
\textbf{phase と own が整合していること}を要求する。

express mode で「phase を書き換えたが own はまだ」という瞬間に
割り込まれると、その瞬間の状態は $\mathrm{wf}(s)$ を満たさない。
型は壊れないが、抽象状態機械との対応が壊れる。
同レポートが「実装との対応は言えない」と述べている箇所が、
まさにここで効いてくる。

\textbf{これが ABCL/1 に atomicity 宣言があった理由である。}
割り込みを許すなら、割り込まれてはいけない区間を宣言できなければならない。
express mode と \texttt{atomic} は分離できない一対の機構である。

\section{メリット}

\subsection{\texttt{where} の利点}

\begin{enumerate}
\item \textbf{同期が宣言的になる。} 有界バッファ、資源プール、
  段階制御（初期化が済むまで受け付けない）などが、
  受理条件としてそのまま書ける。フラグ変数も待ちループも要らない。

\item \textbf{「受理してから条件を見て送り直す」という悪い定石が消える。}
  現在この形を書くと、条件が偽の間メールボックスを自分で埋め続けることになる。
  受信規律の中で待てるようになる。

\item \textbf{受信経路が一本化できる。} 通常 dispatch がガード付き走査に
  なれば、\texttt{select} はその特別な場合になる。
  現在の「\texttt{select} が通常 dispatch に先取りされる」という
  観測可能な不整合が消える。これは表現力を増やす話ではなく、
  後付けで生じた矛盾の解消である。

\item \textbf{既存の機構を再利用できる。} メールボックスの走査は
  \texttt{select} の実装にすでにある。ガード評価を足すだけである。

\item \textbf{型システムをほとんど触らない。} ガードの \ty{bool} 検査が
  増えるだけで、依存型は要らない。
\end{enumerate}

\subsection{express mode の利点}

\begin{enumerate}
\item \textbf{塞がったアクターに手が届く。} 現在 \texttt{now} で待っている
  アクターは石になる。express mode があれば、待っている最中でも
  状態問い合わせ・中止指示・監視に応答できる。
  これは現在の AIPL の最大のモデル的な穴に直接効く。

\item \textbf{デッドロックからの脱出手段になる。} 相互 \texttt{now} で
  固まった二者に、外から \texttt{express send a.abort()} を送れる。
  型では排除できないデッドロックに対する\emph{運用上の}答えになる。

\item \textbf{監視・内省が現実的になる。} \texttt{aios\_*} の内省
  プリミティブは、対象アクターが動いていると意味のある値を返しにくい。
  express mode があれば「いま何をしているか」を当人に聞ける。

\item \textbf{タイムアウトの段階的な扱いが書ける。}
  \texttt{select} の \texttt{timeout} は自分が待つ側の機構だが、
  express mode は\emph{相手を叩く}側の機構である。両方あると
  「一定時間返らなければ問い合わせ、それでも駄目なら中止」が書ける。

\item \textbf{型健全性の証明を壊さない}（\S 5.3）。
\end{enumerate}

\section{デメリットと危険}

\subsection{\texttt{where} の代償}

\begin{center}
\small
\begin{tabular}{@{}lp{8.8cm}@{}}
\toprule
項目 & 内容 \\
\midrule
進行定理が弱まる & 「ガード待ち」という第四の枝が増える（\S 5.2） \\
FIFO でなくなる & メッセージの処理順が到着順と一致しなくなる。
  意味論として明記が必要 \\
飢餓 & あるメッセージのガードが永遠に真にならない場合、
  そのメッセージは永久に残る。公平性は保証できない \\
静的に検出できない誤り & 恒偽のガードを持つメソッドは
  「決して受理されない」が、型検査は通る。
  \texttt{select} の arity 不一致と同じ種類の
  \emph{静かに壊れる}誤りが一つ増える \\
コスト & 配送ごとにメールボックスを走査し、
  各メッセージについてガードを評価する。
  $O(1)$ から $O(n)$ になる \\
ガードの純粋性 & 効果システムが無いので構文的に制限する必要がある（\S 4.4）\\
\bottomrule
\end{tabular}
\end{center}

飢餓と恒偽ガードについては、緩和策がある。
明らかに恒偽なガード（\texttt{false} や定数比較）は警告できる。
また、あるメッセージが一定回数の走査で受理されなかったことを
実行時に記録すれば、診断の手がかりになる。

\subsection{express mode の代償}

\begin{center}
\small
\begin{tabular}{@{}lp{8.8cm}@{}}
\toprule
項目 & 内容 \\
\midrule
\textbf{意味的不変条件が壊れる} & 最大の問題。
  \texttt{atomic} と対で導入しなければ使い物にならない（\S 5.4）\\
前提となる修正 & メソッド呼び出しごとのフレームが必要。
  現在ローカル変数は \texttt{actor.env} に直接置かれ、
  退避・復元されていない（\S 3.3、実測済み）\\
割り込み粒度 & 協調的ポーリングなので文境界でしか割り込めない。
  単一の長い式や組込み関数の内部では割り込めない \\
再入 & express メソッドが通常メソッドを呼ぶ、
  express の中で express が来る、といった入れ子の深さ制限が要る \\
デバッグ可能性 & 実行の非決定性が増える。
  「どこで割り込まれたか」を再現するのが難しい \\
既存プログラムの前提 & 「メソッドは中断されない」を暗黙に仮定している
  既存コードは、express メソッドを一つ足した時点で
  その仮定を失う。既定は非割り込みにすべき \\
証明との対応 & 型健全性は無傷だが、
  抽象状態機械との refinement を主張していた箇所は
  \texttt{atomic} の裏付けが必要になる \\
\bottomrule
\end{tabular}
\end{center}

\subsection{二機構の相互作用}

\begin{itemize}
\item \textbf{express メッセージにガードを許すか。}
  許すと「割り込みたいのに受理されない」という状況が生じ、
  express の目的（緊急の応答）に反する。
  \textbf{express メソッドにはガードを付けられない}とするのが妥当である。

\item \textbf{\texttt{atomic} 区間の中で \texttt{now} を書けるか。}
  書けると、\texttt{atomic} のまま返信を待つことになり、
  その間割り込めない。express mode の目的を損なう。
  \textbf{\texttt{atomic} 区間では \texttt{now} と \texttt{await} を禁止}
  するのが安全である。これは構文的に検査できる。

\item \textbf{ガード評価中に割り込むか。}
  ガードは純粋で短いので、割り込まないのが単純である。
\end{itemize}

\section{段階的な導入案}

いずれも独立に価値があり、この順序なら各段階で退行を測れる。

\begin{center}
\small
\begin{tabular}{@{}rlp{7.2cm}@{}}
\toprule
段階 & 内容 & 備考 \\
\midrule
0 & メソッド呼び出しフレームの導入 &
  ローカル変数を \texttt{actor.env} から出す。
  express mode の前提だが単独でも意味がある潜在バグの修正 \\
1 & \texttt{select} の case に \texttt{where} &
  最小。通常 dispatch の FIFO を変えないので影響が閉じている。
  ガードの \ty{bool} 検査と純粋性検査をここで作る \\
2 & メソッドに \texttt{where}、通常 dispatch を走査化 &
  進行定理に枝が増える。\texttt{select} を統一する好機 \\
3 & \texttt{atomic} 区間 &
  express より\emph{先}に入れる。
  中で \texttt{now} / \texttt{await} を禁止する検査も同時に \\
4 & express mode &
  段階 0 と 3 が前提。既定は非割り込み、
  \texttt{express method} と明示したものだけ割り込める \\
\bottomrule
\end{tabular}
\end{center}

段階 1 と 2 の間、段階 3 と 4 の間で、
既存コーパスの型検査判定と差分テストハーネスを回して測るべきである。
これまでの改修で、見積りが実測と食い違うことが繰り返しあった。

\section{結論}

\subsection{どちらを先に入れるべきか}

\textbf{\texttt{where} を先に入れるべきである。}理由は三つある。

\begin{enumerate}
\item 型システムへの影響が \ty{bool} 検査の追加だけで済む。
\item 失うものが「進行定理に枝が一本増える」と明確に述べられる。
\item \texttt{select} と名前 dispatch の並立という\emph{既存の}不整合を
  解消できる。表現力の追加ではなく、負債の返済としての価値がある。
\end{enumerate}

express mode は、前提となる修正（呼び出しフレーム）と
対で必要な機構（\texttt{atomic}）を抱えており、
しかも既存コードの暗黙の仮定を壊す。
価値は高いが、\texttt{where} のあとに置くのが順序として自然である。

\subsection{全体としてどう評価するか}

この二つを入れると、AIPL は ABCL/1 に確かに近づく。
そして\textbf{その代償は型システムではなく、推論の難しさとして現れる}。

型健全性は無傷である。保存は影響を受けず、
逐次性は形式モデルが最初から落としているので express mode も証明を壊さない。
進行定理だけが \texttt{where} によって枝を一本増やす。

しかし、プログラムについて\emph{個別に}証明していた性質
---- 哲学者の食事問題の資源デッドロック自由のような ----
は、割り込みが入ると抽象状態機械との対応を \texttt{atomic} で
裏づけ直す必要がある。

先の比較レポートで
「失った表現力が型付けを容易にしている」と書いた。本稿はその裏返しである。
\textbf{表現力を戻すと、型は保てるが、個別の証明の手間が増える}。
どちらを重く見るかが判断の分かれ目であり、
研究目的が「実装された言語の健全性を機械検証する」ことにあるなら、
\texttt{where} までは進み、express mode は \texttt{atomic} の設計が
固まってから、というのが妥当な線である。

\part{二つの目的に向けた言語仕様の考察}
\label{part:purpose}

\section{二つの目的とその要求}

ここから視点を変える。AIPL が次の二つを目的に持つとしたら、
どのような言語仕様であるべきか。

\begin{itemize}
\item \textbf{分散 AI エージェント言語} ---- 複数ノードに分散した
  エージェントが、モデル呼び出し・記憶・タスク・プロトコルを通じて協調する。
  実装にはすでに \texttt{ai\_call}、\texttt{model\_generate}、
  \texttt{aios\_memory\_*}、\texttt{aios\_task\_*}、\texttt{protocol\_*}、
  \texttt{remote(...)}、HTTP ゲートウェイが揃っている。
\item \textbf{実時間 OS（フィジカル AI）の記述言語} ---- センサ・アクチュエータを
  持つ実機の上で、周期処理と割り込みを記述する。
  別リポジトリで Xinu 上の実行系（プリエンプティブ優先度、tickless タイマ）と
  AIPL からのコード生成が進んでいることを前提とする。
\end{itemize}

\subsection{分散 AI エージェント言語としての要求}

\begin{enumerate}
\item \textbf{待ち時間が長く、しかも予測できない。}
  モデル呼び出しは秒単位で、失敗もする。
  \texttt{now} で待つと、その間アクターは何もできない。
\item \textbf{失敗が普通である。}タイムアウト、部分的な成功、リトライが日常であり、
  「返るか返らないか」ではなく「何が返り、返らなかったときどうするか」を
  書けなければならない。
\item \textbf{結果が構造を持つ。}モデルの出力は JSON 相当の構造である。
  現在の AIPL の型は \ty{int} / \ty{float} / \ty{string} / \ty{bool} /
  \ty{unit} / 配列 / アクターだけで、\textbf{レコードも直和も無い}。
  構造をすべて文字列で運ぶことになる。
\item \textbf{どこへ何が出ていくかが重要である。}機外へ出る呼び出しと
  ローカル推論は区別されなければならない。
  監査・秘匿・課金のいずれの観点からも必要である。
\item \textbf{対話の順序に規律がある。}\texttt{protocol\_*} は
  すでにセッションプロトコルを実行時に検査しているが、型ではない。
\end{enumerate}

\subsection{実時間 OS（フィジカル AI）としての要求}

\begin{enumerate}
\item \textbf{時間が仕様である。}周期、期限、最悪実行時間、ジッタが
  正しさの一部になる。「いつか終わる」では不十分である。
\item \textbf{有界であること。}メモリ、ループ回数、メールボックス長、
  再帰深さ ---- すべてに上限が要る。
  動的割り付けは初期化後には禁止したい。
\item \textbf{割り込みが本質である。}センサ・タイマ・通信の到着は割り込みであり、
  第 \ref{part:mech} 部の express mode がまさにこれである。
  そして割り込まれてはいけない区間（\texttt{atomic}）が必ず要る。
\item \textbf{ブロックが致命的である。}\texttt{now} で無期限に待つ制御ループは
  期限を守れない。
\item \textbf{数値の決定性。}浮動小数はハードウェアによって時間も結果も揺れる。
  制御ループでは整数か固定小数点が望ましい。
  なお AIPL では \texttt{/} が \ty{int} 同士でも \ty{float} を返す仕様である。
\item \textbf{物理量の取り違えが事故になる。}\texttt{m/s} と
  \texttt{rad/s} を足しても現在の型システムは何も言わない。
\end{enumerate}

\subsection{両者は多くの点で正反対である}

\begin{center}
\small
\begin{tabular}{@{}lll@{}}
\toprule
 & 分散 AI エージェント & 実時間 OS \\
\midrule
待ち時間        & 秒単位、予測不能 & マイクロ秒〜ミリ秒、上界必須 \\
失敗            & 日常。設計に織り込む & 許されない（設計で排除する） \\
データ          & 構造化・可変長・スキーマ推測 & 固定長・既知 \\
メモリ          & 動的でよい & 初期化後は静的 \\
非決定性        & 本質的（モデル出力） & 排除したい \\
\texttt{now}    & 使いたいが長く待つ & 期限なしでは使えない \\
浮動小数        & 問題ない & 避けたい \\
外部通信        & 中心的な機能 & 限定・帯域保証つき \\
優先度          & 概念が薄い & 中心的 \\
\bottomrule
\end{tabular}
\end{center}

同じ言語で両方を\emph{無差別に}許すと、どちらの保証も得られない。
実時間側は「この部分は割り付けもブロックもしない」と言えなくなり、
エージェント側は不要な制約を負う。

\section{提案 ---- 一つの中核と二つのプロファイル}

\subsection{先例 ---- Ada の Ravenscar プロファイル}

この形には直接の先例がある。Ada は汎用言語だが、
高信頼・実時間用途に対して\textbf{Ravenscar プロファイル}という
制限された部分集合を定義した。動的なタスク生成を禁じ、
待ち合わせの形を限定し、解析可能性を確保する。
言語を分けるのではなく、\emph{同じ言語の中に検証可能な部分集合を切る}
という設計である。

同様に、分散側には E 言語という先例がある。E は\textbf{ブロックを一切許さず}、
eventual send と promise だけで書かせる。デッドロックを言語から
排除するための選択である。

AIPL が採るべきはこの二つを合わせた形である。

\begin{center}
\small
\begin{tabular}{@{}lp{5.0cm}p{5.0cm}@{}}
\toprule
 & プロファイル A（\texttt{agent}） & プロファイル R（\texttt{realtime}） \\
\midrule
方針 & ブロックを避け、失敗を型に持つ &
  有界性と期限を静的に保証する \\
許す effect & \texttt{net}, \texttt{ai}, \texttt{fs}, \texttt{mem},
  \texttt{time} & \texttt{io}（宣言した装置のみ）, \texttt{time} \\
禁止 & 無期限 \texttt{now} & \texttt{now}（期限なし）、動的割り付け、
  非有界ループ、再帰、浮動小数 \\
必須 & 失敗の扱い、リモートインタフェース宣言 &
  周期／期限／優先度の宣言、\texttt{atomic} \\
\bottomrule
\end{tabular}
\end{center}

\subsection{すでにある土台 ---- capability 分類}

好都合なことに、実装にはすでに capability の分類が入っている。
プリミティブごとに分類が付けられ、実行時に列挙できる。
実測すると 17 分類あった。

\begin{lstlisting}[language={}]
Core.Math  Core.Array  Core.Introspection  Console  Time
Actor  Actor.Introspection
AIOS.Kernel  AIOS.Service  AIOS.Task  AIOS.Memory  AIOS.Event
AIOS.Model  AIOS.Model.OpenAI  AIOS.Model.Gemini  AIOS.Remote
Protocol.Session  UI.SDL  Web
\end{lstlisting}

\begin{finding}
この分類は\textbf{実行時メタデータにとどまっており、型検査器は一切使っていない}。
\texttt{capabilities()} と \texttt{capability\_prims(c)} で内省できるだけである。
\end{finding}

つまりプロファイル分離のための\emph{データはすでにある}。
足りないのは、これを静的な effect として\emph{メソッドに伝播させ、
宣言と照合する}機構だけである。姉妹実装（Python 版 AIPL）では
効果注釈が gradual に検査されており、方式の前例も内部にある。

\subsection{effect を静的にする}

メソッドに effect を宣言でき、本体が使うプリミティブと送信先の effect の
和が宣言に含まれることを検査する。

\begin{lstlisting}[language=AIPL]
class Planner {
  // 機外へ出ることを宣言している
  method plan(goal) : string !{ai, net} {
    reply(ai_call("plan: " ++ goal));
  }
}

class Controller {
  // io と time だけ。net も ai も使えない
  method step(sensor) : int !{io, time} {
    reply(clamp(sensor));
  }
}
\end{lstlisting}

これで得られるものは大きい。

\begin{itemize}
\item \textbf{機外に出ないローカル推論を型で示せる。}
  たとえばオンデバイス推論のプリミティブに \texttt{\{ai\}} だけを与え、
  \texttt{net} を含めなければ、「AI を使うが機外へは出ない」ことが
  シグネチャに現れる。
\item \textbf{実時間側で禁止すべき effect を列挙できる。}
  \texttt{net} や \texttt{ai} を持つメソッドは
  プロファイル R では呼べない、と機械的に言える。
\item \textbf{監査ができる。}どのアクターがどの外部資源に触るかが
  静的に集計できる。
\end{itemize}

\section{中核に必要な仕様}

両プロファイルに共通して必要なものを挙げる。
いずれも現在の穴を埋める作業でもある。

\begin{enumerate}
\item \textbf{型の穴を閉じる。}
  \texttt{sender} が \ty{any} で送信が無検査、リモート送信先が \ty{any}、
  アクター型が単一化で区別されない、\texttt{become} が置換後の
  インタフェースを検査しない ---- これらは分散でも実時間でも致命的である。
  分散では相手を間違え、実時間では存在しないメソッドへ送っても
  実行時まで分からない。

\item \textbf{リモートインタフェース記述。}
  対向ノードの \texttt{class C \{ method m(x) : T \}} だけを書いた
  記述を読み込み、\texttt{remote(...)} に型を付ける。
  戻り値型注釈を必須にしたのはこの前提づくりであった。
  実時間メッシュでもエージェント連携でも、これが無いと境界が \ty{any} のままである。

\item \textbf{レコードと直和。}
  現在は構造化データを表現できない。
  エージェント側は JSON 相当、実時間側はセンサフレームやコマンドを
  構造として運びたい。直和があれば「成功か失敗か」も型で表せる。

\item \textbf{\texttt{where} ガード}（第 \ref{part:mech} 部）。
  エージェント側では「必要なデータが揃うまで受理しない」、
  実時間側では「状態機械の段階に合わない指令を受理しない」に効く。
  どちらも本質的である。

\item \textbf{\texttt{==} と \texttt{!=}。}
  現在これらは字句解析器と評価器にあるのに\emph{文法規則が無く書けない}。
  状態機械を書く言語としては欠落が大きい。
\end{enumerate}

\section{プロファイル A ---- 分散 AI エージェント}

\subsection{無期限 \texttt{now} を禁止し、期限付きにする}

エージェントにとって最悪の事態は、モデル呼び出しの返りを待って
アクターが石になることである。\texttt{now} を廃するのではなく、
\textbf{期限と失敗時の分岐を必須にする}。

\begin{lstlisting}[language=AIPL]
var r = now planner.plan(goal) timeout 5000 else "fallback";
\end{lstlisting}

型としては \texttt{now ... timeout ... else e} の全体が
$\rho$ 型の式であり、\texttt{else} 節も同じ型を持つことを要求すればよい。
既存の \texttt{select} の \texttt{timeout} と同じ機構で実装できる。

これは進行定理に対して\textbf{良い方向}に働く。現在の進行定理の第三の枝
「\texttt{await} 待ち」は、期限付き \texttt{now} だけを使うプログラムでは
有限時間で解消するので、その範囲でデッドロック自由が言える。
言い換えれば、\textbf{期限を必須にすることは型で保証できる進行性を増やす}。

\subsection{失敗を型に持つ}

期限切れ、モデルの拒否、ネットワーク断は日常である。
直和があれば戻り値型として書ける。

\begin{lstlisting}[language=AIPL]
method plan(goal) : Result<string, Error> !{ai, net} { ... }
\end{lstlisting}

現在は「失敗したら文字列に \texttt{"ERROR: ..."} を入れる」しかなく、
呼び出し側が確認を忘れても型検査は通る。

\subsection{セッションプロトコルを型へ上げる}

\texttt{protocol\_*} は実行時にプロトコル状態を検査している。
これを型レベルのセッション型へ上げることは、健全性レポートが
すでに次の課題として挙げている。
戻り値型注釈を必須にしたことで、$!m(T).?R.\cdots$ の $R$ が
インタフェースに書かれている状態になった。前提は整っている。

\subsection{構造化出力のスキーマ検査}

モデル出力を型付きの値へ取り込む箇所は、必ず動的検査になる。
そこを\textbf{一点に閉じ込める}のが型システムの役目である。

\begin{lstlisting}[language=AIPL]
// 検査に失敗したら Err になる。以降は型が保証される
var p = parse<Plan>(now planner.plan(goal) timeout 5000 else "");
\end{lstlisting}

\section{プロファイル R ---- 実時間 OS／フィジカル AI}

\subsection{時間を宣言する}

周期、期限、優先度をメソッドの宣言に置く。

\begin{lstlisting}[language=AIPL]
class Balancer {
  // 1kHz、期限 800us、優先度 10
  periodic method tick() : unit !{io, time}
      period 1000us deadline 800us priority 10 {
    var s = read_gyro();
    set_motor(pid(s));
  }
}
\end{lstlisting}

これは型システムの仕事ではなく\emph{スケジューリング可能性解析}の仕事だが、
言語が宣言を持てば解析器を書ける。Xinu 上の実行系が
プリエンプティブ優先度と tickless タイマを持つなら、
宣言をそのままスケジューラの設定へ落とせる。

\subsection{有界性を静的に強制する}

プロファイル R では次を禁止または制限する。

\begin{center}
\small
\begin{tabular}{@{}ll@{}}
\toprule
対象 & 扱い \\
\midrule
\texttt{new}（動的アクター生成） & 初期化フェーズのみ許可 \\
配列の伸長（\texttt{array\_push}） & 禁止。容量は宣言する \\
\texttt{while} & 禁止。上界の明らかな \texttt{for} だけを許す \\
再帰・自分あて送信の連鎖 & 禁止（呼び出しグラフで検査） \\
期限なし \texttt{now} / \texttt{await} & 禁止 \\
浮動小数 & 禁止または固定小数点に置換 \\
\texttt{net}, \texttt{ai}, \texttt{fs} effect & 禁止 \\
\bottomrule
\end{tabular}
\end{center}

これらはいずれも\textbf{構文と effect の検査で機械的に実施できる}。
依存型は要らない。
\texttt{while} の禁止は既存の \texttt{reply} 線形性検査
（ループ本体の \texttt{reply} を「2 回以上」と判定する走査）と
同じ形の解析で足りる。

\subsection{数値 ---- 固定小数点と単位}

浮動小数を避けるなら、\ty{int} を基本にした固定小数点型が要る。
ここで、本セッションで整理した数値の扱いが直接効いてくる。
現在 \texttt{/} は \ty{int} 同士でも \ty{float} を返す。
プロファイル R ではこれを許せないので、
\texttt{div} と \texttt{mod}、あるいは固定小数点の \texttt{/} を
別に用意することになる。

さらにフィジカル AI では\textbf{物理単位}を型に載せる価値が大きい。

\begin{lstlisting}[language=AIPL]
var v : float<m/s>   = 1.2;
var w : float<rad/s> = 0.3;
var x = v + w;          // 単位不一致として型エラーにしたい
\end{lstlisting}

単位は型の添字として扱えるので、依存型は不要である。
乗除で指数を足し引きする規則を単一化に入れるだけで済み、
制御ソフトの事故の実際的な原因を一つ潰せる。

\subsection{割り込みとしての express mode}

第 \ref{part:mech} 部の express mode は、実時間側では選択肢ではなく\textbf{必然}である。
センサ・タイマ・通信の到着は割り込みであり、
\texttt{express method} がハンドラ、\texttt{atomic} が
クリティカルセクションに対応する。

そして第 \ref{part:mech} 部で挙げた前提と制約が、そのまま実時間の作法になる。

\begin{itemize}
\item メソッド呼び出しフレームの導入（割り込みハンドラは入れ子で走る）。
\item \texttt{atomic} 区間では \texttt{now} / \texttt{await} を禁止。
\item ハンドラ自身も有界でなければならない。
  \textbf{express メソッドはプロファイル R の制約を必ず満たす}
  ---- ループ禁止、割り付け禁止、期限つき ---- とすべきである。
\end{itemize}

\subsection{有界メールボックスと溢れ方針}

無限に伸びるキューは実時間では使えない。
容量と、溢れたときの方針を宣言させる。

\begin{lstlisting}[language=AIPL]
class Sensor mailbox 8 overflow drop_oldest { ... }
\end{lstlisting}

溢れが起きうることは型では消せないが、
\emph{どうするかを宣言させる}ことで、暗黙の失敗を無くせる。

\section{二つのプロファイルの境界}

分散エージェントと実時間制御は同じシステムの中で協調する
---- エージェントが判断し、実時間側が実行する。
両者をつなぐ規律が要る。

\subsection{方向の規律}

\begin{center}
\small
\begin{tabular}{@{}llp{6.0cm}@{}}
\toprule
方向 & 許す送信 & 理由 \\
\midrule
A $\to$ R & \texttt{send} のみ &
  実時間側の応答を待つとエージェント側が塞がる。
  結果が要るなら R 側から通知させる \\
R $\to$ A & 期限付き \texttt{future} のみ &
  実時間側は決してブロックしてはならない。
  \texttt{now} は使えない \\
\bottomrule
\end{tabular}
\end{center}

すなわち\textbf{実時間側からエージェント側への同期呼び出しを型で禁止する}。
これは effect と期限の検査で機械的に実施できる。
「制御ループが LLM の返りを待って期限を落とす」という
最も起こりやすい事故を、言語が構造的に防ぐ。

\subsection{境界アクター}

境界はアクターとして明示する。

\begin{lstlisting}[language=AIPL]
// エージェント側の判断を受け取り、実時間側へ渡すだけの境界
class Bridge profile boundary {
  var target = 0;

  // A 側からは send のみ
  express method set_target(v) : unit !{} { target = v; }

  // R 側は即座に読める
  method get_target() : int !{} { reply(target); }
}
\end{lstlisting}

境界アクターは\textbf{状態の受け渡しだけを行い、計算しない}。
express で書き込み、通常メソッドで読む。
これなら実時間側は待たない。

\section{型健全性と証明への影響}

提案した仕様が、既存の形式的保証にどう影響するかを整理する。

\begin{center}
\small
\begin{tabular}{@{}lp{8.6cm}@{}}
\toprule
仕様 & 影響 \\
\midrule
effect 系 & 保存・進行に影響なし。effect は型の\emph{添え物}として
  独立に定式化できる（型と effect の対で保存を述べる標準的な形） \\
レコード・直和 & 標準的な拡張。保存の証明が場合分けで長くなるだけ \\
期限付き \texttt{now} & \textbf{進行定理を強める}。
  期限付き待ちは有限時間で解消するので、
  第三の枝が「有限時間後に解消する待ち」に格上げできる \\
\texttt{where} & 進行定理に「ガード待ち」の枝を追加（第 \ref{part:mech} 部）\\
express + \texttt{atomic} & 型健全性は無傷。
  プログラム個別の不変条件は \texttt{atomic} で裏づけ直す必要がある \\
有界性の制約 & 型システムの外側の構文・グラフ解析。
  健全性の証明には影響しないが、\emph{停止性}という
  現在言えていない性質をプロファイル R では言えるようになる \\
物理単位 & 型の添字として単一化に入れるだけ。依存型は不要 \\
\bottomrule
\end{tabular}
\end{center}

注目すべきは、\textbf{プロファイル R では現在言えていない性質が言えるようになる}
ことである。健全性レポートは「停止性は言えない ---- メソッドが自分あてに
送れば無限に走る」と述べているが、プロファイル R は
まさにそれを禁止するので、停止性と応答時間の上界が主張できる。

\section{実装の順序}

これまでの改修で、見積りが実測と食い違うことが繰り返しあった。
以下は順序の提案であり、各段階で既存コーパスの判定と
差分テストハーネスを回して測るべきである。

\begin{center}
\small
\begin{tabular}{@{}rlp{6.6cm}@{}}
\toprule
 & 内容 & 位置づけ \\
\midrule
1 & \texttt{==} / \texttt{!=} の文法追加 & 数行。欠落の解消 \\
2 & 型の穴を閉じる（\texttt{sender}、アクター型の区別、\texttt{become}）&
  両プロファイルの前提 \\
3 & effect 注釈と検査（既存 capability 分類を静的化）&
  \textbf{プロファイル分離の土台}。ここが要 \\
4 & リモートインタフェース記述 & 分散の境界を閉じる \\
5 & レコードと直和 & エージェント側の構造化データ、失敗の型 \\
6 & 期限付き \texttt{now} & 進行性が強まる。A の必須要件 \\
7 & \texttt{where} ガードと受信経路の統一 & 第 \ref{part:mech} 部 \\
8 & 呼び出しフレーム、\texttt{atomic}、express mode &
  第 \ref{part:mech} 部。R の必須要件 \\
9 & 有界性の制約とプロファイル宣言 & R の閉じ込め \\
10 & 周期・期限・優先度の宣言とスケジューリング解析 &
  Xinu 実行系との接続 \\
11 & 物理単位 & あれば事故が減る。優先度は低い \\
12 & セッション型 & 長期目標 \\
\bottomrule
\end{tabular}
\end{center}

段階 3 が要である。effect が静的になれば、
プロファイルの区別も、境界の方向の規律も、
実時間側の禁止事項も、すべて同じ機構の上に載る。
そして既存の capability 分類がそのまま使えるので、
分類を新たに設計する必要はない。

\section{まとめ ---- どのような言語仕様とすべきか}

\begin{enumerate}
\item \textbf{一つの中核と二つのプロファイル}にする。
  二つの目的の要求は正反対であり、無差別に許すとどちらの保証も得られない。
  Ada の Ravenscar プロファイルという先例がある。

\item \textbf{プロファイルを分ける機構は effect/capability 系}である。
  実装にはすでに 17 種の capability 分類があり、
  実行時メタデータにとどまっている。これを静的にすることが最優先である。

\item \textbf{中核では型の穴を閉じ、構造化データを入れる。}
  \texttt{sender} と \texttt{remote} の \ty{any}、
  アクター型の非区別、\texttt{become} の未検査、
  レコード・直和の不在、\texttt{==} の欠落 ---- これらは
  分散でも実時間でも支障になる。

\item \textbf{エージェント側は「ブロックしない」を型で強制する。}
  期限付き \texttt{now} を必須にし、失敗を直和で表す。
  これは進行定理を\emph{強める}方向の変更である。

\item \textbf{実時間側は「有界である」を静的に強制する。}
  ループ・割り付け・再帰・浮動小数・外部通信を制限し、
  周期と期限を宣言する。express mode と \texttt{atomic} は
  選択肢ではなく必然である。
  そしてこのプロファイルでは、現在言えていない停止性と
  応答時間の上界が主張できるようになる。

\item \textbf{境界は方向を規律する。}
  実時間側からエージェント側への同期呼び出しを型で禁止する。
  「制御ループがモデルの返りを待って期限を落とす」という
  最も起こりやすい事故を、言語が構造的に防ぐ。
\end{enumerate}

第 \ref{part:mech} 部で「表現力を戻すと、型は保てるが個別の証明の手間が増える」と書いた。
第 \ref{part:purpose} 部の結論はその続きである。
\textbf{目的を二つに分けて宣言させれば、それぞれの目的に見合った保証を
別々に主張できる。}一つの言語で一つの保証を目指すより、
一つの中核で二つの保証を目指すほうが、この二つの目的には適している。


\part{AIPL 理想仕様（案）}
\label{part:spec}

第 \ref{part:purpose} 部の考察を、仕様として明文化する。
以下は「こうあるべき」という案であり、現状の実装とは異なる。
現状との差分は \S\ref{sec:delta} にまとめる。

\section{設計原則}

\begin{enumerate}
\item[\textbf{P1}] \textbf{アクターが唯一の並行単位である。}
  共有変数を持たない。状態はアクターに閉じ、
  外からはメッセージでしか触れない。

\item[\textbf{P2}] \textbf{通信は三種、そして待つなら期限を持つ。}
  \texttt{send}（待たない）、\texttt{now}（待つ）、
  \texttt{future}（後で受け取る）。
  \textbf{無期限に待つ構文は存在しない。}

\item[\textbf{P3}] \textbf{返信はちょうど一度、宣言された型で行う。}
  戻り値型は必ず宣言する。推論は補助であって省略の許可ではない。

\item[\textbf{P4}] \textbf{受理は名前とガードだけで決まり、受信経路は一本である。}
  同じメッセージが二つの経路で競合しない。

\item[\textbf{P5}] \textbf{効果は宣言し、検査する。}
  プロファイルは効果と構文の制限として定義される。
  新しい概念を足すのではなく、既にある capability 分類を静的にする。

\item[\textbf{P6}] \textbf{境界はインタフェース宣言でのみ型付ける。}
  \ty{any} は仕様に存在しない。
  リモートも外部公開も、宣言された interface を通す。

\item[\textbf{P7}] \textbf{実時間プロファイルでは有界性を静的に強制する。}
  ループ、割り付け、再帰、待ち、外部通信のすべてに上限か禁止を課す。
\end{enumerate}

\section{構文}

\subsection{宣言}

\begin{lstlisting}[language={}]
program    ::= decl*
decl       ::= class | interface | record | variant | global

class      ::= "class" ID profile? mailbox? "{" field* method* "}"
profile    ::= "profile" ("agent" | "realtime" | "boundary")
mailbox    ::= "mailbox" INT overflow?
overflow   ::= "overflow" ("drop_oldest" | "drop_newest" | "error")

interface  ::= "interface" ID "{" sig* "}"
sig        ::= "method" ID "(" params ")" ":" type effects? ";"

record     ::= "record" ID "{" ID ":" type ("," ID ":" type)* "}"
variant    ::= "variant" ID "{" ID ("(" type ")")?
                          ("|" ID ("(" type ")")?)* "}"

field      ::= "var" ID ":" type "=" expr ";"
method     ::= "express"? "method" ID "(" params ")" ":" type
               effects? guard? timing? block
params     ::= (ID ":" type ("," ID ":" type)*)?
effects    ::= "!" "{" effect ("," effect)* "}"
guard      ::= "where" expr
timing     ::= ("period" dur)? ("deadline" dur)? ("priority" INT)?
dur        ::= INT ("us" | "ms" | "s")
global     ::= stmt
\end{lstlisting}

現状との主な違いは、\textbf{仮引数とフィールドに型を書く}こと、
\textbf{戻り値型が必須}であること、
そして \texttt{interface} / \texttt{record} / \texttt{variant} /
\texttt{profile} / \texttt{mailbox} / \texttt{express} /
\texttt{where} / \texttt{timing} が加わることである。

\subsection{型}

\begin{lstlisting}[language={}]
type ::= "int" | "bool" | "string" | "unit"
       | "float" | "float" "<" unit ">"        // 物理単位つき
       | "fix" "<" INT "," INT ">"             // 固定小数点（整数部,小数部ビット）
       | type "[" INT "]"                      // 固定長配列
       | type "[]"                             // 可変長配列（プロファイル A のみ）
       | "future" "<" type ">"
       | ID                                    // actor / interface / record / variant
       | ID "<" type ("," type)* ">"           // 型構築子の適用
unit ::= ID ("^" INT)? (("*"|"/") ID ("^" INT)?)*   // m, m/s, rad/s, kg*m/s^2
\end{lstlisting}

\textbf{\ty{any} は無い。}
現状 \ty{any} が使われている三箇所は次のように置き換える。

\begin{center}
\small
\begin{tabular}{@{}lll@{}}
\toprule
現状 & 理想仕様 \\
\midrule
\texttt{sender} : \ty{any} & \texttt{sender} : \texttt{Reply<}$\rho$\texttt{>}
  （返信のみ可能な能力） \\
リモート送信先 : \ty{any} & 宣言された \texttt{interface} 型 \\
未注釈メソッドの戻り値 : \ty{any} & 戻り値型は必須なので存在しない \\
\bottomrule
\end{tabular}
\end{center}

\subsection{式}

\begin{lstlisting}[language={}]
expr ::= INT | FLOAT | STRING | "true" | "false" | ID
       | expr binop expr | "!" expr
       | expr "." ID                                  // レコードのフィールド
       | ID "{" ID "=" expr ("," ID "=" expr)* "}"    // レコード生成
       | ID "(" args ")"                              // 関数・構築子
       | "new" ID "(" args ")"                        // プロファイル A / 初期化のみ
       | "now" target "." ID "(" args ")" "timeout" dur "else" expr
       | "express" "now" target "." ID "(" args ")" "timeout" dur "else" expr
       | "future" target "." ID "(" args ")" "deadline" dur
       | "await" expr "timeout" dur "else" expr
       | "match" expr "{" ("case" pat "->" expr)+ "}"
       | "(" expr ")"
binop ::= "+" | "-" | "*" | "/" | "div" | "mod" | "++"
        | "==" | "!=" | "<" | ">" | "<=" | ">=" | "&&" | "||"
target ::= ID | "self" | "remote" "(" STRING "," STRING ")" ":" ID
\end{lstlisting}

要点を四つ挙げる。

\begin{itemize}
\item \textbf{\texttt{now} と \texttt{await} は期限と \texttt{else} を伴う。}
  期限なしの形は文法に存在しない（\textbf{P2}）。
\item \texttt{==}、\texttt{!=}、\texttt{\&\&}、\texttt{||}、\texttt{!} を入れる。
  現状これらは書けない。
\item \texttt{/} は浮動小数の除算、\texttt{div} / \texttt{mod} は整数除算とする。
  \ty{int} 同士の \texttt{/} が \ty{float} を返す現状の曖昧さを解消する。
\item リモート送信先は \texttt{remote("host:port","actor") : IFace} と
  interface を明記する。これで境界の \ty{any} が消える。
\end{itemize}

\subsection{文}

\begin{lstlisting}[language={}]
stmt  ::= "var" ID ":" type "=" expr ";"
        | ID "=" expr ";"
        | expr ";"
        | "reply" "(" expr ")" ";"
        | "send" target "." ID "(" args ")" ";"
        | "express" "send" target "." ID "(" args ")" ";"
        | "if" "(" expr ")" stmt ("else" stmt)?
        | "for" ID "in" expr ".." expr "do" stmt          // 有界
        | "while" "(" expr ")" stmt                       // プロファイル A のみ
        | "atomic" block
        | "select" "{" case+ timeout? "}"
        | "become" ID "(" args ")" ";"
        | "match" expr "{" ("case" pat "->" stmt)+ "}"
        | block
case  ::= "case" ID "(" params ")" guard? "->" block
\end{lstlisting}

\texttt{for} は上下限が明らかな有界ループで、両プロファイルで使える。
\texttt{while} はプロファイル A のみとする。

\section{型システム}

\subsection{判断の形}

型付けは\textbf{型と効果の対}で述べる。

\[
\Gamma \vdash_{C.m} e : \tau\ !\ \varepsilon
\]

$\Gamma$ は変数環境、$C.m$ は検査中のメソッド（\texttt{reply} の帰属先）、
$\tau$ が型、$\varepsilon$ が効果集合である。

\subsection{主要な規則}

\begin{mathpar}
\inferrule*[right=\textsc{Reply}]
  {\mathit{ret}(C,m) = \tau \\ \Gamma \vdash_{C.m} e : \tau\ !\ \varepsilon}
  {\Gamma \vdash_{C.m} \texttt{reply}(e) : \ty{unit}\ !\ \varepsilon}
\and
\inferrule*[right=\textsc{Send}]
  {\Gamma \vdash a : \ty{actor}(C') \\
   C'.m' : (\vec{\tau}) \to \tau'\ !\ \varepsilon' \\
   \Gamma \vdash e_i : \tau_i\ !\ \varepsilon_i}
  {\Gamma \vdash \texttt{send}\ a.m'(\vec{e}) : \ty{unit}\ !\ \textstyle\bigcup\varepsilon_i}
\and
\inferrule*[right=\textsc{Now}]
  {\Gamma \vdash a : \ty{actor}(C') \\
   C'.m' : (\vec{\tau}) \to \tau'\ !\ \varepsilon' \\
   \Gamma \vdash e_i : \tau_i\ !\ \varepsilon_i \\
   \Gamma \vdash e_{\mathit{alt}} : \tau'\ !\ \varepsilon_a}
  {\Gamma \vdash \texttt{now}\ a.m'(\vec{e})\ \texttt{timeout}\ d\
    \texttt{else}\ e_{\mathit{alt}} : \tau'\ !\
    \varepsilon' \cup \varepsilon_a \cup \textstyle\bigcup\varepsilon_i}
\and
\inferrule*[right=\textsc{Method}]
  {\Gamma, \overline{x{:}\tau} \vdash g : \ty{bool}\ !\ \emptyset \\
   \Gamma, \overline{x{:}\tau} \vdash_{C.m} \mathit{body} : \ty{unit}\ !\ \varepsilon_b \\
   \varepsilon_b \subseteq \varepsilon \\
   \mathit{replies}(\mathit{body}) = 1}
  {\Gamma \vdash \texttt{method}\ m(\overline{x{:}\tau}) : \tau_r\ !\varepsilon\
    \texttt{where}\ g\ \{\mathit{body}\}}
\end{mathpar}

三点が重要である。

\begin{enumerate}
\item \textbf{\textsc{Send} は callee の効果 $\varepsilon'$ を伝播させない。}
  効果は別アクターで起こるからである。
\item \textbf{\textsc{Now} は伝播させる。}呼び出し側が待つので、
  待っている間の性質は呼び出し側の責任になる。
  \textbf{この一点だけで境界の方向の規律が導かれる}
  ---- プロファイル R は \texttt{net} や \texttt{ai} を持てないので、
  それらを持つメソッドを \texttt{now} で呼べない。
\item \textbf{ガードは効果を持てない}（$\varepsilon = \emptyset$）。
  受理の試行で状態が変わってはならない。
\end{enumerate}

\subsection{返信の規律}

$\mathit{replies}(\mathit{body}) = 1$ は「全実行パスでちょうど一度
\texttt{reply} する」である。現在は上界を構文で、
下界を注釈時にのみ検査しているが、
戻り値型が必須になれば両方を常に課せる。

\texttt{select} の case 本体と express メソッドの本体は
\emph{別の}返信文脈である。それぞれ対応するメッセージの
$\mathit{ret}$ に照合される。

\subsection{物理単位}

単位は型の添字であり、乗除で指数を足し引きする。

\begin{mathpar}
\inferrule*{ }{\ty{float}\langle u_1\rangle \times \ty{float}\langle u_2\rangle
  \to \ty{float}\langle u_1 u_2\rangle}
\and
\inferrule*{ }{\ty{float}\langle u\rangle + \ty{float}\langle u\rangle
  \to \ty{float}\langle u\rangle}
\end{mathpar}

加減は単位が一致していなければ型エラーとする。依存型は不要である。

\section{効果}

\begin{center}
\small
\begin{tabular}{@{}lll@{}}
\toprule
効果 & 意味 & 対応する既存 capability \\
\midrule
\texttt{mut}   & 自分の状態を書き換える & --- \\
\texttt{time}  & 時刻の取得、待機 & \texttt{Time} \\
\texttt{io}    & 宣言した装置への入出力 & \texttt{UI.SDL}、GPIO 等 \\
\texttt{mem}   & 動的割り付け（\texttt{new}、配列の伸長）&
  \texttt{Core.Array}、\texttt{Actor} \\
\texttt{net}   & ノード外への通信 &
  \texttt{Web}、\texttt{AIOS.Remote} \\
\texttt{ai}    & モデル推論 &
  \texttt{AIOS.Model}（機外へ出るものは \texttt{net} も併記）\\
\texttt{fs}    & 永続化 & \texttt{AIOS.Memory} \\
\texttt{log}   & 観測のみの出力 & \texttt{Console}、\texttt{AIOS.Event} \\
\bottomrule
\end{tabular}
\end{center}

対応表があるので、既存の 17 分類から機械的に初期値を与えられる。
新しく設計する必要はない。

\begin{remark}
\texttt{ai} と \texttt{net} を分けるのが要点である。
オンデバイス推論は \texttt{\{ai\}} だけを持ち \texttt{net} を持たない。
これにより「AI を使うが機外へは出ない」がシグネチャに現れる。
フィジカル AI では、この区別が安全性と秘匿の両方で意味を持つ。
\end{remark}

\section{プロファイル}

\begin{center}
\small
\begin{tabular}{@{}lp{3.9cm}p{3.9cm}p{2.4cm}@{}}
\toprule
 & \texttt{agent} & \texttt{realtime} & \texttt{boundary} \\
\midrule
許す効果 &
  すべて &
  \texttt{mut}, \texttt{time}, \texttt{io}, \texttt{log} &
  \texttt{mut} のみ \\
\texttt{new} &
  可 & 初期化フェーズのみ & 不可 \\
配列 &
  可変長可 & 固定長のみ & 固定長のみ \\
\texttt{while} &
  可 & 不可（\texttt{for} のみ）& 不可 \\
再帰・自己送信の連鎖 &
  可 & 不可 & 不可 \\
\texttt{now} / \texttt{await} &
  期限必須 & 期限必須。かつ相手の効果が
  許容集合に収まる場合のみ & 不可 \\
浮動小数 &
  可 & \texttt{fix<i,f>} または
  単位つき \ty{float} のみ & 不可 \\
\texttt{timing} 宣言 &
  任意 & \textbf{必須} & --- \\
\texttt{mailbox} 宣言 &
  任意 & \textbf{必須} & 必須 \\
\texttt{express} &
  可 & 可（ハンドラも R の制約に従う）& 可 \\
\bottomrule
\end{tabular}
\end{center}

プロファイルは\textbf{型システムの外側の構文・グラフ解析}で強制する。
依存型は要らない。そして \textsc{Now} の効果伝播規則があるので、
「実時間アクターがエージェントを同期で呼ぶ」は
\emph{効果の検査だけで}禁止される。

\section{実行モデル}

\subsection{アクターの状態}

各アクターは通常キューと express キューを持つ。

\begin{center}
\small
\begin{tabular}{@{}ll@{}}
\toprule
状態 & 内容 \\
\midrule
\textsf{idle}    & 受理可能なメッセージを待っている \\
\textsf{running} & メソッド本体を実行中 \\
\textsf{blocked} & 期限付きで返信を待っている（期限で必ず解ける）\\
\textsf{atomic}  & \texttt{atomic} 区間の実行中。express を保留する \\
\bottomrule
\end{tabular}
\end{center}

ABCL/1 の dormant / active / waiting に対応する。
違いは \textsf{blocked} が必ず有限時間で解けることである（\textbf{P2}）。

\subsection{受理規則 ---- 一本の経路}

\begin{enumerate}
\item express キューに、名前が一致し\emph{ガードを持たない}
  express メソッドに対応するメッセージがあれば、それを受理する。
  \textsf{atomic} 中は保留する。
\item そうでなければ、通常キューを到着順に走査し、
  \textbf{名前が一致し、かつガードが真}である最初のメッセージを受理する。
\item どちらも無ければ \textsf{idle} で待つ。
  新しい到着、または状態の変化があれば再走査する。
\end{enumerate}

\texttt{select} は「対象を指定した名前集合に限り、期限を付けた」
規則 2 の特別な場合である。これで受信経路が一本になり、
現在の先取り競合（\S 3 の \texttt{select} の注意）が消える。

\subsection{\texttt{atomic}}

\begin{itemize}
\item \texttt{atomic} 区間では express を受理しない。
\item \texttt{atomic} 区間では \texttt{now} / \texttt{await} を書けない
  （構文検査）。待ちながら割り込みを止めることを許さない。
\item \texttt{atomic} 区間は有界でなければならない
  （プロファイル R では \texttt{for} のみ）。
\end{itemize}

\subsection{返信}

\texttt{reply} は、いま受理しているメッセージに一度だけ答える。
express ハンドラや \texttt{select} の case で受理したメッセージは
それぞれ独立の返信先を持つ。

\section{現状との差分}
\label{sec:delta}

\subsection{足すもの}

\begin{center}
\small
\begin{tabular}{@{}ll@{}}
\toprule
項目 & 効く目的 \\
\midrule
\texttt{==} \texttt{!=} \texttt{\&\&} \texttt{||} \texttt{!} & 両方 \\
\texttt{record} / \texttt{variant} と \texttt{match} & 主にエージェント \\
\texttt{interface} 宣言とリモートの型付け & 両方 \\
効果注釈と検査 & 両方（プロファイルの土台）\\
\texttt{profile} 宣言 & 両方 \\
\texttt{where} ガード & 両方 \\
\texttt{express} と \texttt{atomic} & 主に実時間 \\
\texttt{mailbox} 容量と溢れ方針 & 主に実時間 \\
\texttt{timing}（周期・期限・優先度）& 実時間 \\
\texttt{fix<i,f>} と物理単位 & 実時間 \\
\texttt{for} 有界ループ & 実時間 \\
メソッド呼び出しフレーム（実装）& express の前提 \\
\bottomrule
\end{tabular}
\end{center}

\subsection{変えるもの}

\begin{center}
\small
\begin{tabular}{@{}lll@{}}
\toprule
 & 現状 & 理想仕様 \\
\midrule
戻り値型 & 省略可（推論）& \textbf{必須}。推論は補助 \\
仮引数・フィールド & 型を書かない & 型を書く \\
\texttt{now} / \texttt{await} & 期限なし & \textbf{期限と \texttt{else} 必須} \\
\texttt{/} & \ty{int} 同士でも \ty{float} &
  \texttt{/} は浮動小数、\texttt{div}/\texttt{mod} が整数 \\
\texttt{sender} & \ty{any}、送信は無検査 &
  \texttt{Reply<}$\rho$\texttt{>}。返信のみ \\
リモート送信先 & \ty{any} & \texttt{interface} 型 \\
受信経路 & 名前 dispatch と \texttt{select} が並立 &
  一本（\texttt{select} は特別な場合）\\
\texttt{become} & 引数のみ検査 &
  インタフェースの一致を検査 \\
アクター型 & 単一化で区別されない & クラス名で区別 \\
\bottomrule
\end{tabular}
\end{center}

\subsection{禁止するもの（プロファイル R）}

期限なしの待ち、\texttt{while}、再帰と自己送信の連鎖、
初期化後の \texttt{new} と配列の伸長、浮動小数（単位つきと固定小数点を除く）、
\texttt{net} / \texttt{ai} / \texttt{fs} 効果、
\texttt{atomic} 中の待ち。

\section{この仕様で言えるようになること}

\begin{center}
\small
\begin{tabular}{@{}lll@{}}
\toprule
性質 & 現状 & 理想仕様 \\
\midrule
型健全性（保存・進行）& 言える & 言える（効果つきに拡張）\\
返信がちょうど一度 & 注釈時のみ & \textbf{常に言える} \\
境界の型 & \ty{any} & \textbf{言える}（interface）\\
デッドロック自由 & 一般には言えない &
  \textbf{期限付き待ちのみなので、有限時間で解ける} \\
停止性 & 言えない & \textbf{プロファイル R では言える} \\
応答時間の上界 & 言えない &
  \textbf{プロファイル R では解析できる}（宣言があるので）\\
機外へ出ないこと & 言えない & \textbf{効果で言える} \\
単位の整合 & 言えない & \textbf{言える} \\
公平性 & 言えない & 言えない（\texttt{where} で悪化する）\\
\bottomrule
\end{tabular}
\end{center}

現状「言えない」とされていた四つ ---- デッドロック自由、停止性、
応答時間、外部流出 ---- が、\textbf{プロファイルを宣言させることで
言えるようになる}。これが第 \ref{part:purpose} 部の結論の具体的な中身である。

引き換えに公平性は悪化する。\texttt{where} を入れると
ガードが永遠に真にならないメッセージが残りうるからで、
これは仕様として明記し、診断で補うべき部分である。

\section{最小の到達点}

全部を一度に作る必要はない。
次の四つが揃った時点で「二つの目的に耐える言語」と言える。

\begin{enumerate}
\item \textbf{効果注釈と検査}（既存 capability 分類の静的化）。
  プロファイル分離、境界の方向の規律、外部流出の保証がここに乗る。
\item \textbf{期限付き \texttt{now} / \texttt{await} のみ}。
  デッドロック自由が言えるようになる。
\item \textbf{\texttt{interface} 宣言}。\ty{any} が消える。
\item \textbf{\texttt{profile realtime} の有界性検査}。
  停止性と応答時間が言えるようになる。
\end{enumerate}

\texttt{record} / \texttt{variant}、物理単位、\texttt{where}、
express mode は、この四つの上に順に足していける。


\begin{thebibliography}{9}

\bibitem{abcl1}
Akinori Yonezawa, Jean-Pierre Briot, Etsuya Shibayama.
\newblock Object-Oriented Concurrent Programming in ABCL/1.
\newblock \emph{OOPSLA 1986}, pp.~258--268.

\bibitem{abclbook}
Akinori Yonezawa (ed.).
\newblock \emph{ABCL: An Object-Oriented Concurrent System}. MIT Press, 1990.

\bibitem{cmp}
本リポジトリ \texttt{docs/aipl\_vs\_abcl1\_ja.pdf} ---
  AIPL と ABCL/1 の機能・モデル比較。

\bibitem{snd}
本リポジトリ \texttt{docs/aipl\_soundness\_ja.pdf} ---
  AIPL$^-$ の型健全性・型安全性の Coq 検証、哲学者の食事問題、
  および「保証しないこと」の一覧。

\bibitem{rep}
本リポジトリ \texttt{docs/aipl\_reply\_inference\_ja.pdf} ---
  \texttt{reply} からの戻り値型推論と \texttt{select} における
  返信の帰属の扱い。

\bibitem{guide}
本リポジトリ \texttt{docs/aipl\_guide\_ja.pdf} --- 現状の言語ガイド。

\bibitem{ravenscar}
Alan Burns, Brian Dobbing, Tullio Vardanega.
\newblock Guide for the use of the Ada Ravenscar Profile in high integrity
systems.
\newblock \emph{Ada Letters} XXIV(2), 2004.
\newblock （同じ言語の中に検証可能な部分集合を切るという設計の先例）

\bibitem{elang}
Mark Miller, E. Dean Tribble, Jonathan Shapiro.
\newblock Concurrency Among Strangers: Programming in E as Plan Coordination.
\newblock \emph{TGC 2005}.
\newblock （ブロックを一切許さず eventual send と promise だけで書かせる設計）

\bibitem{lustre}
Nicolas Halbwachs, Paul Caspi, Pascal Raymond, Daniel Pilaud.
\newblock The synchronous data flow programming language LUSTRE.
\newblock \emph{Proceedings of the IEEE} 79(9), 1991.
\newblock （実時間記述における有界性と時間の宣言）

\end{thebibliography}

\end{document}
